% sec:rk-geodesic-application — Application: Geodesic Integration
\section{Application: Geodesic Integration}
\label{sec:rk-geodesic-application}

The preceding sections developed the Runge--Kutta integrator as an
abstract machine: feed it a right-hand side $f$, a state vector $y$,
and a step size $h$, and it advances the solution while controlling the
local error.  For the ray tracer, the abstraction collapses to a
specific system --- Hamilton's equations for a photon in curved
spacetime --- and the solver parameters acquire direct physical
meaning.  Each piece of the abstract machinery --- state vector,
right-hand side, error norm, step-size controller --- acquires a
concrete counterpart in the geodesic codebase.

\paragraph{The state vector.}

The abstract state vector $y$ of \cref{eq:ivp} becomes the
eight-component phase-space state of the Hamiltonian formulation
(\cref{sec:geo-formulations}):
\coderef{Geodesic.Types}{HamiltonianState}
%
\begin{equation}
  y = (x^\mu,\, p_\mu) \,, \qquad \mu = 0, 1, 2, 3 \,,
  \label{eq:geodesic-state}
\end{equation}
%
comprising four position coordinates and four covariant momenta.  The
independent variable $t$ in the abstract ODE is the affine parameter
$\lambda$ along the geodesic.  For null rays, the affine rescaling
freedom (\cref{sec:geo-null}) lets the ray tracer set $E = -p_t = 1$
at the initial point without loss of generality, so the affine
parameter carries no independent physical scale --- it is simply a
label for how far along the ray the integrator has advanced.  The eight
numbers in $y$ encode everything the integrator needs: the photon's
location in spacetime and the direction it is travelling.

\paragraph{The right-hand side.}

The abstract function $f(t, y)$ in \cref{eq:ivp} becomes the
Hamiltonian right-hand side derived in \cref{sec:geo-formulations}.
\coderef{Geodesic.Hamiltonian}{hamiltonianRHS}
Writing $f = (\dot{x}^\mu, \dot{p}_\mu)$, the two components are
%
\begin{equation}
  \begin{aligned}
  \dot{x}^\mu &= g^{\mu\nu}\,p_\nu \,, \\
  \dot{p}_\mu &= \tfrac{1}{2}\,(\pd{\mu}\,g_{\alpha\beta})\,
    \dot{x}^\alpha\,\dot{x}^\beta \,.
  \end{aligned}
  \label{eq:geodesic-rhs}
\end{equation}
%
The velocity equation lifts covariant momenta to contravariant
velocities via the inverse metric; the momentum equation updates
$p_\mu$ using the metric's spatial variation along the trajectory.
The velocity $\dot{x}^\alpha = g^{\alpha\beta}\,p_\beta$ from the
first equation feeds into the second, so the codebase evaluates the
inverse metric before the momentum update.
\implnote{The function \texttt{makeRHS} in
\codemodule{Geodesic.Solver} wraps \texttt{hamiltonianRHS} by
evaluating the inverse metric at the current position before each
call, forming a closure of type
\texttt{HamiltonianState -> StateVec}.}

Each evaluation of $f$ requires the inverse metric $g^{\mu\nu}$ at
the current position and the 40 independent metric derivatives
$\pd{\mu}\,g_{\alpha\beta}$.  The codebase obtains these derivatives
by passing a polymorphic metric function through forward-mode automatic
differentiation (\cref{ch:autodiff}), so adding a new spacetime
requires only the metric --- no hand-coded Christoffel symbols.  The
inverse metric is evaluated once per RHS call and reused in both the
velocity computation and the momentum contraction.

\paragraph{From abstract stages to physical evaluations.}

Each Dormand--Prince step (\cref{sec:rk-dormand-prince}) computes
seven stage slopes $k_1, \ldots, k_7$, where each $k_i$ is an
eight-component vector of the same dimension as $y$.  In the geodesic
context, each stage slope is a pair
$k_i = (\Delta x^\mu_i,\, \Delta p_{\mu,i})$, and each stage
evaluation requires a full metric evaluation: the inverse metric at
the stage point, followed by the 40 metric derivatives via AD.  The
FSAL property (\cref{sec:rk-fsal}) saves one of these evaluations per
accepted step, reducing the per-step cost from seven to six metric
evaluations.

Over a typical ray trace of $10^4$ accepted steps, the integrator
performs roughly $6 \times 10^4$ RHS evaluations --- each involving an
inverse metric lookup, an AD metric-derivative evaluation, and the Hamilton's
equations contraction.  For a Kerr--Schild metric, the inverse metric
is algebraically simple ($\eta^{\mu\nu} - f\,l^\mu l^\nu$), so the AD
evaluation dominates.  The adaptive controller concentrates these
evaluations where the dynamics demand them: a typical ray might
complete the weak-field approach phase in a few dozen large steps
but require several hundred small steps during a close pass around
the photon sphere.

\paragraph{The error norm and phase-space scaling.}

The adaptive controller (\cref{sec:rk-adaptive}) applies the scaled
error norm $\mathrm{err}$ of \cref{eq:rk-error-norm} to all eight
components of the phase-space state simultaneously.  For the geodesic
system this means that position errors and momentum errors compete
for the same tolerance budget.  Whether this is sensible depends on
the coordinate system.

In Kerr--Schild Cartesian coordinates with $M = 1$, typical
coordinates satisfy $\abs{x^\mu} \sim 1$--$10^2$ (the observer sits
at $r_0 \sim 50M$--$500M$) and typical covariant momenta satisfy
$\abs{p_\mu} \sim 1$ (with the $E = 1$ normalisation).  With the
default $\epsilon_a = \epsilon_r = 10^{-10}$, a position component at
$r \sim 100M$ has tolerance scale
$\epsilon_a + \epsilon_r \cdot 100 \approx 10^{-8}$ (the $\epsilon_a$
term is negligible), while a momentum component of order one has
tolerance scale $\epsilon_a + \epsilon_r \cdot 1 = 2 \times 10^{-10}$.
The absolute tolerance $\epsilon_a$ also protects components passing
through zero --- for instance, $p_z$ passes through zero for a photon
crossing the equatorial plane.  A
uniform tolerance pair $(\epsilon_a, \epsilon_r)$ therefore provides
sensible scaling without requiring separate position and momentum
tolerances, provided the coordinate system keeps both sets of
components within a few orders of magnitude of each other.
\warnnote{In coordinate systems where position and momentum
components differ by many orders of magnitude --- for example,
spherical coordinates with $r \sim 10^6 M$ and
$p_\varphi \sim 10^{-3}$ --- separate tolerances or a rescaled
state vector would be more appropriate.  The codebase avoids this
issue by working exclusively in Cartesian Kerr--Schild coordinates.}

\paragraph{The integration pipeline.}

The function \texttt{integrateGeodesic} assembles the RHS closure,
adaptive controller, and termination logic into a single ray-tracing
call:
%
\begin{enumerate}
\item Construct the RHS closure from the polymorphic metric function
  and inverse metric lookup.
\item Compute the initial stage $k_1^{(0)} = f(\lambda_0, y_0)$ from
  the initial state.
\item Enter the adaptive loop: at each step, call
  \texttt{dormandPrinceStep} with the current state and $k_1$
  (from FSAL), compute the error norm, accept or reject, adjust the
  step size, and check termination conditions
  (\cref{sec:geo-termination}).
\item Return the final state paired with the termination reason ---
  horizon capture, celestial-sphere escape, disc intersection, or
  step-count exhaustion.
\end{enumerate}
%
The termination check occurs \emph{before} each Dormand--Prince step,
not after: if the ray has already crossed the horizon or escaped,
there is no point computing six metric evaluations only to discard
the result.  The disc-crossing check, which requires interpolation
between the current and proposed states, occurs after an accepted
step.
\implnote{The diagnostic variant \texttt{traceGeodesic} in
\codemodule{Geodesic.Solver} records the full trajectory ---
$(\lambda, \text{state}, \mathcal{H})$ at every accepted step ---
enabling the conservation studies of \cref{sec:geo-conservation}.}

\paragraph{Connecting tolerances to image quality.}

The solver tolerances $\epsilon_a$ and $\epsilon_r$ ultimately
determine the visual quality of the rendered black hole image.  A ray
whose accumulated error exceeds the angular resolution of a single
pixel produces a displaced pixel --- invisible in isolation, but
visible as noise or banding in regions of high image gradient, such as
the edge of the shadow or the inner rim of the accretion disc.

The codebase's default tolerances $\epsilon_a = \epsilon_r = 10^{-10}$
(\cref{eq:rk-defaults}) are conservative: the Hamiltonian constraint
violation $\abs{\mathcal{H}}$ tracks $\epsilon_a$
(\cref{eq:H-violation-bound}), and the conservation monitoring of
\cref{sec:geo-conservation} confirms that position errors remain
negligible compared with the pixel angular resolution --- of order
$10^{-4}$~rad per pixel for a typical $1024 \times 1024$ image.  For preview renders, relaxing to
$\epsilon_a = \epsilon_r = 10^{-6}$ reduces the step count (and
therefore the render time) while keeping visible artefacts below the
perceptible threshold at typical display resolutions.  The tolerances
thus function as a quality--performance dial: tighter for
publication-quality images, looser for interactive exploration.
\xrefnote{The camera model that maps the final ray direction to a
pixel coordinate is developed in \cref{ch:ray-tracing}; the tolerance
choice interacts with the angular resolution set by the image
dimensions.}
